\documentclass[12pt, a4paper, UTF8, openany]{ctexbook}

\usepackage{geometry}%设置页边距
\geometry{left=1.54cm, right=1.54cm, top=2.18cm, bottom=2.18cm}

\usepackage{tikz}%绘图
\usetikzlibrary{calc}%计算

\usepackage{circuitikz}%画电路图

\usepackage{amsmath, amssymb}%数学

\usepackage{array}%表格基础设置
\usepackage{longtable}%跨页表格
\usepackage{booktabs}%表格水平线

\usepackage[colorlinks = true, 
            linkcolor = black]{hyperref}%超链接跳转

\usepackage{color}

\linespread{1.05}%行间距
\setlength{\parskip}{8pt}%段落间距
\setlength{\parindent}{0pt}%不缩进

\newcommand{\lxor}{\oplus}
\newcommand{\bi}{\leftrightarrow} % biconditional
\newcommand{\bigand}{\bigwedge}
\newcommand{\bigor}{\bigvee}
\newcommand{\true}{\mathbf{T}}
\newcommand{\false}{\mathbf{F}}

\newcommand{\field}[1]{\mathbb{#1}}%域

% \setcounter{chapter}{-1}

\title{Discrete Mathematics}
\author{1001}

\begin{document}
\maketitle

\tableofcontents % use LaTeXmk
\cleardoublepage

\chapter{Logic and Proofs}
{\large $\mathbf{The\ Only\ Truth\ is\ the\ Truth\ Table.}$}

    \section{Basic}
    \textbf{Proposition}: A declarative sentence that is either True or False, but not both. \\
    not declarative: $x + 1 = 3$ \\
    not T or F: Paradox. e.g. This sentence is false. \\
    allow its truth value be undetermined: He will go to Paris.

    \textbf{Proposition Varible}: $p, q, r$, etc.

    \textbf{Atomic Proposition}: the simplest proposition.

    Given proposition $p \to q$:
    \begin{longtable}[h]{*{2}{l}}
        \textbf{Converse} 逆命题 & $q \to p$ \\
        \textbf{Inverse} 否命题 & $\lnot p \to \lnot q$ \\
        \textbf{Contrapositive} 逆否命题 & $\lnot q \to \lnot p$
    \end{longtable}

    $p$: \textbf{Hypothesis} / \textbf{Antecedent} / \textbf{Premise} \\
    $q$: \textbf{Conclusion} / \textbf{Consequence}

    \textbf{Tautology}: always True; \\
    \textbf{Contradiction}: always False; \\
    \textbf{Contingency}: sometimes True, sometimes False; \\
    \textbf{Satisfiable}: 存在赋值方式使为真, Tautology + Contingency

    Precedence: $\forall/\exists \;>\; \lnot \;>\; \land/\lor \;>\; \to/\bi$

    \textbf{Theorem}: a somewhat important statement. \\
    \textbf{Proposition} / Facts / Results: less important. \\
    \textbf{Axiom} 公理 / \textbf{Postulate} 假设: can be used as premises. \\
    \textbf{Lemma}: a less important theorem that is helpful for other results. \\
    \textbf{Corollary} 推论: a theorem that can be established directly from a theorem. \\
    \textbf{Conjecture} 猜想: guess to be true.

    \section{Logical Operators}
        \subsection{Defines}
        \begin{longtable}[h]{l | l | l | l | l}
            Negation 否定 & $\lnot x$ / $!x$ / $\sim x$ & not & & $1 - x$ \\
            Conjunction 合取 & $x \land y$ & and & & $xy$ \\
            Disjunction 析取 & $x \lor y$ & or & & $x + y - xy$ \\
            Exclusive Or & $x \lxor y$ & xor & & \\
            Conditional / Implication 蕴涵 & $x \to y$ & if, then & $\lnot x \lor y$ & $1 - x + xy$ \\
            Biconditional / Equivalence 等价 & $x \bi y$ & iff & $\lnot (x \lxor y)$ & \\
        \end{longtable}

        x only if y: $x \to y$ \\
        x unless y: $y \to \lnot x$ \\
        x, y both occur: $x \land y$, no causality. \\
        avoid using \textit{any}, 因为自然语言下同时有 任意一个/某一个 的含义。

        \begin{circuitikz} %[scale = 0.8, transform shape]
            \draw (0, 0) node[not port] {};
            \node at (0, -1) {Inverter};
            \draw (5, 0) node[and port] {};
            \node at (4.5, -1) {And Gate};
            \draw (10, 0) node[or port] {};
            \node at (9.5, -1) {Or Gate};
        \end{circuitikz}

        \subsection{Rules}
        Logical Equivalence: 
        \begin{itemize}
            \item Identity: $x \land \true \equiv x$, $x \lor \false \equiv x$
            \item Domination: $x \land \false \equiv \false$, $x \lor \true \equiv \true$
            \item Idempotent: $x \land/\lor x \equiv x$
            \item Double Negation: $\lnot(\lnot x) \equiv x$
            \item Communtative \& Associative Laws: $\land/\lor$ 分别满足交换 \& 结合律
            \item Distribution: $\land/\lor$ 分别对 $\lor/\land$ 分配 \\
                $\begin{cases}
                    x \lor (y \land z) & \equiv (x \lor y) \land (x \lor z) \\
                    x \land (y \lor z) & \equiv (x \land y) \lor (x \land z)
                \end{cases}$
            \item De Morgan's Laws: flip operator \\
                $\begin{cases}
                    \lnot (x \land y) & \equiv \lnot x \lor \lnot y \\
                    \lnot (x \lor y) & \equiv \lnot x \land \lnot y
                \end{cases}$
            \item Absorption Laws: \\
                $\begin{cases}
                    x \lor (x \land y) & \equiv x \\
                    x \land (x \lor y) & \equiv x
                \end{cases}$
        \end{itemize}

        \subsection{Normal Forms}
        Simple Conjunctive/Disjunctive Formula: only consists of $\land$/$\lor$ and $\lnot$

        Conjunctive/Disjunctive Normal Form (CNF / DNF): $\bigand(\lor)$ / $\bigor(\land)$ \\
        主合取/析取范式：外层为合取/析取。CNF / DNF 必存在，构造性证明，略。

        CNF / DNF 不唯一，故引出 \textbf{Principle/Major C/D NF} (P*NF/M*NF)（主合/析取范式），强制约定每个基本项包含所有 $n$ 个变元。
        P*NF 是唯一的。

        极小项：$\bigand_{i = 1}^{n}$, $m_{0} \sim m_{2^n - 1}$；\\
        极大项：$\bigor_{i = 1}^{n}$, $M_{0} \sim M_{2^n - 1}$；\\
        编号：from left to right, $i$-th bit $= [p_i]$。\\
        e.g. $(p_1 \land \lnot p_2 \land \lnot p_3) = m_{100} = m{4}$; $(\lnot p_1 \lor p_2) = M_{01} = M_{1}$

        Construct P*NF base on *NF: for each simple formula with less than $n$ varibles, 
        adjust, and merge the same formulas \\
        a. $A \; \equiv \; (A \land p_k) \lor (A \land \lnot p_k)$ \\
        b. $A \; \equiv \; (A \lor p_k) \land (A \lor \lnot p_k)$

        转换 PCNF/PDNF，其中 $\sim$ 为按位取反：
        $\mathop{\bigand}\limits_{t \in S} M_{t} = \mathop{\bigor}\limits_{t \in U \setminus S} m_{\sim t}$；
        $\mathop{\bigor}\limits_{t \in S} m_{t} = \mathop{\bigand}\limits_{t \in U \setminus S} M_{\sim t}$ \\
        证明：PDNF $\to$ PCNF：每个 $t \in S$ 表示一种成真赋值，则每个 $t \in U \setminus S$ 表示一种成假赋值。
        为使结果为真，则每种成假赋值均需有一项不成立，即补集的按位取反。易知该操作可逆，则 PCNF $\to$ PDNF 亦得。

    \section{Predicates and Quantifiers}
        \subsection{Defines}
        \textbf{Predicate} 谓词: refer to a property, propositonal function $P$, $P(x_1, \cdots x_n)$

        \textbf{Quantifier} 量词: the extent to which $P$ is true over a range of elements \\
        a. universal quantifier 全称量词: $\forall$ \\
        b. existential quantifier 存在量词: $\exists$ \\
        c. uniqueness quantifier 唯一量词: $\exists!$/$\exists_1$

        \textbf{Domain}: the domain of $x$; the domain consists of $a, b, c, \cdots$ \\
        \textbf{Scope}: the scope of $\forall/\exists$, etc.

        Bound Occurance 约束出现: $x$ is restricted by quantifiers. \\
        Free Occurance 自由出现: $x$ is not restricted by any quantifiers. \\
        Closed Form 闭式: no free varibles, all varibles is in some scope. 

        \subsection{Rules}
        \begin{itemize}
            \item Negation: \\
                a. $\lnot \forall \; \equiv \; \exists \lnot$ \\
                b. $\lnot \exists \; \equiv \; \forall \lnot$
            \item Distribution: turn $\to/\bi$ to $\land/\lor$ first \\
                a. under the same restriction: only the following two are correct
                \begin{itemize}
                    \item $\forall x(A(x) \land B(x)) \; \equiv \; \forall x A(x) \land \forall x B(x)$
                    \item $\exists x(A(x) \lor B(x)) \; \equiv \; \exists x A(x) \lor \exists x B(x)$
                \end{itemize}
                b. under different restriction: 
                    $\forall/\exists x(A(x) \land/\lor B) \; \equiv \; \forall/\exists x A(x) \land/\lor B$
        \end{itemize}

        \subsection{Prenex Normal Form}
        前束范式: all quantifier are at the beginning, and their scopes last to the end. \\
        e.g. $\forall x \exists y \forall z (P(x, y) \lor Q(x, y, z) \lor \lnot R(z))$

        PNF 必存在，构造性证明，略。\\
        PNF 不唯一。

    \section{Rules of Inference}
    a \textbf{valid argument}: $(\bigand\text{premises}) \to \text{conclusion}$

    \textbf{fallacy}: 谬误

        \subsection{Rules}
        \begin{longtable}[h]{l | l | l}
            Modus Ponens & 肯定前件 & 
            $\begin{cases}
                x \\ x \to y    
            \end{cases} \implies y$ \\
            \addlinespace

            Modus Tollens & 否定后件 & 
            $\begin{cases}
                \lnot y \\ x \to y
            \end{cases} \implies \lnot x$ \\
            \midrule

            Hypothetical Syllogism & 假言三段论 & 
            $\begin{cases}
                x \to y \\ y \to z
            \end{cases} \implies x \to z$ \\
            \addlinespace

            Disjunctive Syllogism & 析取三段论 & 
            $\begin{cases}
                x \lor y \\ \lnot x
            \end{cases} \implies y$ \\
            \midrule

            Addition & & $x \implies (x \lor y)$ \\
            \addlinespace

            Simplification & & $(x \land y) \implies x$ \\
            \addlinespace

            Conjunction & & 
            $\begin{cases}
            x \\ y    
            \end{cases} \implies (x \land y)$ \\
            \midrule

            Resolution & 消解 & 
            $\begin{cases}
                (x \lor y) \\ (\lnot x \lor z)
            \end{cases} \implies (y \lor z)$ \\
            \addlinespace

            Constructive Dilemma & &
            $\begin{cases}
                (x \to y) \land (u \to v) \\ x \lor u
            \end{cases} \implies (y \lor v)$ \\
            \addlinespace

            Destructive Dilemma & &
            $\begin{cases}
                (x \to y) \land (u \to v) \\ (\lnot y \lor \lnot v)
            \end{cases} \implies (\lnot x \lor \lnot u)$ \\
            \midrule

            Universal Instantiation (UI) & & $\forall xP(x) \implies P(u)$ \\
            \addlinespace
            Existential Instantiation (EI) & & $\exists xP(x) \implies P(u) \text{ for some u}$ \\
            \addlinespace
            Universal Generalization (UG) & & $P(u) \text{ for an arbitrary } u \implies \forall xP(x)$ \\
            \addlinespace
            Existential Generalization (EG) & & $P(u) \text{ for some } u \implies \exists xP(x)$ \\
            \midrule

            Universal Modus Ponens & & 
            $\begin{cases}
                \forall x(P(x) \to Q(x)) \\ P(u)
            \end{cases} \implies Q(u)$ \\
            \addlinespace

            Universal Modus Tollens & &
            $\begin{cases}
                \forall x(P(x) \to Q(x)) \\ \lnot Q(u)
            \end{cases} \implies \lnot P(u)$
        \end{longtable}

        prep.: using / from

        e.g.
        \begin{longtable}[h]{l l l}
            1. & $\lnot t$ & premise 3. \\
            2. & $r \to t$ & premise 2. \\
            3. & $\lnot r$ & modus tollens using 1, 2 \\
            4. & $\lnot r \lor \lnot s$ & addition from 3 \\
            5. & $\lnot (r \lor s)$ & De Morgan's
        \end{longtable}

        \subsection{Fallacies}
        \begin{itemize}
            \item fallacy of affirming the conclusion: \\
                $\begin{cases}
                    x \to y \\ y
                \end{cases} \nRightarrow  x$
            \item fallacy of denying the hypothesis: \\
                $\begin{cases}
                    x \to y \\ \lnot x
                \end{cases} \nRightarrow  \lnot y$
        \end{itemize}

    \section{Proof Methods}
    Direct Proof.

    Conditional Proof (CP) 附加前提证明: to prove $A \implies (B \to C)$, use $B$ as a premise. 
    try to prove $(A \land B) \to C$.

    Proof by Contraposition: use $\lnot q \to \lnot p$ to prove $p \to q$.

    Proof by Contradiction 反证

    Proof by Counterexample 找反例: use $\lnot P(u)$ to prove $\lnot \forall xP(x)$.

    Existence Proof: \\
    a. Constructive: find a witness, that $P(u)$ is true. \\
    b. Nonconstructive

    Exhaustive Proof / Proof by Exhaustion 穷举 \\
    Proof by Cases 分讨

    Tips:
    \begin{itemize}
        \item Vacuous Proof: if $p$ is false, then $p \to q$ is immediately true.
        \item Trivial Proof: if $q$ is true, then $p \to q$ is always true.
        \item proof of equivalence: \\
            a. to prove $p \bi q$, proving $(p \to q) \land (q \to p)$. \\
            b. to prove $p_1 \bi p_2 \bi \cdots \bi p_n$, proving $\bigand (p_i \to p_{i \bmod n + 1})$.
        \item \textbf{WLOG}: without loss of generality.
    \end{itemize}

\chapter{Relations}
Can't use emptyset for {\heiti Cartesian Product}.

    \section{Definitions}
    A binary relation $R$ from $A$ to $B$ is a subset of $A \times B$: $R \subset A \times B$

    If $(a, b) \in R$, then $a$ is related to $b$. The \textbf{order} matters. Denoted as $a R b$ or $(a, b) \in R$. 
    Otherwise, $a \, R\!\!\!\!/ \; b$ or $(a, b) \notin R$.

    The length of tuple is called {\heiti degree}. e.g. The degree of $A \times \cdots \times A$ is $n$.\\
    $A_i$ is called the domains of $R$.

    \subsection{Representation}
        \textbf{Order!}

        \begin{itemize}
            \item \textbf{directed} bipartite graph
            \item \textbf{directed} graph for $R \subset A \times A$
            \item 01-matrix: $R: A \to B$, $M_{i, j} = [(a_i, b_j) \in R]$ \\
                $M_{R \cap S} = M_R \land M_S$; $M_{R \cup S} = M_R \lor M_S$ \\
                $M_{S \circ R} = M_R \odot M_S$ a.k.a. $M_{i, j} = \lor (M_{S_{i, k}} \land M_{R_{k, j}})$ {\heiti 矩阵顺序不一样！}\\
                $M_{R^n} = M_R^n$
        \end{itemize}

    \subsection{Special Relations}
        \begin{itemize}
            \item 空关系 $\emptyset$
            \item 恒等关系 $I_A = \{(a, a): \forall a\}$
            \item 全域关系 $U_A = \{(a, b): \forall a \forall b\}$
        \end{itemize}

    \subsection{Properties}
        \begin{itemize}
            \item 自反性 Reflexive: $\forall x \in A, x R x$ \\
                Each node has a self-loop. \\
                All elements on the main diagonal are 1.
            \item 反自反性 Irreflexive: $\forall x \in A, (x, x) \notin R$ \\
                Irreflexive is \textbf{not} the negation of reflexive ! \\
                $R$ can be neither reflexive nor irreflexive. $R$ is both reflexive and irreflexive iff $A = \emptyset$ (however, Cartesian Product asks $A \neq \emptyset$ in $A \times A$, so?). \\
                Each node must not have a self-loop. \\
                All elements on the main diagonal are 0.
            \item 对称性 Symmetric: $R = R^{-1}$ (a.k.a. $\forall x, y \in A \; (x, y) \in R \iff (y, x) \in R$) \\
                （等价地）All edges are bi-directed. \\
                All are same across the main diagonal.
            \item 反对称性 Antisymmetric: $(x, y) \in R \land (y, x) \in R \; \implies \; x = y$ \\
                Antisymmetric is \textbf{not} "not symmetric" ! \\
                For any two different vertices, there are at most one directed edge link them. \\
                All 1's are across from 0's.
            \item 非对称性 Asymmetric: $(x, y) \in R \; \implies \; (y, x) \notin R$ \\
                Completely no symmetric. Asymmetric $\iff$ Antisymmetric $\land$ Irreflexive. \\
                No self-loops, and for any two different vertices, there are at most one directed edge link them. \\
                All elements on the main diagonal are 0, and all 1's are across from 0's.
            \item 传递性 Transitive: $(x, y), (y, z) \in R \; \implies \; (x, z) \in R$
        \end{itemize}
        
    \subsection{Operations}
        Inverse Relation:
        $R^{-1} := \{(b, a): (a, b) \in R\}$.\\
        e.g. $R = \{(a, b): a < b\}$, then $R^{-1} = \{(a, b): a > b\}$.\\
        Note that if $R = \{(a, b): p(a, b)\}$, $R^{-1}$ is \textbf{not} $\{(a, b): \lnot p(a, b)\}$ !

        Properties of Inverse Relation:\\
        a. $(R^{-1})^{-1} = R$ \\
        b. $(R \; \cap/\cup/- \; S)^{-1} = R^{-1} \;\cap/\cup/- \; S^{-1}$ \\
        c. $R \subset S \iff R^{-1} \subset S^{-1}$ \\
        d. $(R \circ S)^{-1} = S^{-1} \circ R^{-1}$

        Composition:
        $S \circ R: A \mathop\to\limits^{R} B \mathop\to\limits^{S} C$
        
        Power:\\
        a. $R^0 := I$\\
        b. $R^n := R^{n - 1} \circ R$\\
        c. $R^{-n} := (R^{-1})^n$
        
        \textbf{Thm}. $R$ is transitive iff $R^n \subset R$ for all $n \in \field{N}_+$
        
        \textbf{Thm}. Composition of relations satisfies the associative law.
        
    \section{Closures of Relations}
    \subsection{definitions}
        For any property P, the "P closure" of $R$, if it exists, is defined as the "smallest" subset of $A \times A$ containing $R$ with P.

        \begin{itemize}
            \item contains $R$
            \item has property $P$
            \item "smallest": for every $R \subset T$ s.t. $T$ has property P, $S \subset T$.\\
                因此，闭包是合法集合，且是所有合法 $T$ 的交集。若把集合的包含关系看作合法集之间的偏序关系，则存在闭包说明存在{\heiti 唯一}的极小元，这对 P 本身有所限制。
        \end{itemize}

        closures:
        \begin{itemize}
            \item reflexive closure: $r(R) = R \cup I_A$
            \item symmetric closure: $s(R) = R \cup R^{-1}$
            \item transitive closure: $t(R) = R^\ast = \mathop\cup\limits_{n \in \field{N}_+} R^n$
        \end{itemize}

    \subsection{Connectivity Relation}
        a.k.a. transitive closure

        $(a, b) \in R^n$: there exists a path from $a$ to $b$ with length $n$

        $|A| = n$, if there is a path $a \rightsquigarrow b$, then there must exist a path with length $\leqslant n$. 
        Therefore, $t(R) = R \cup \cdots \cup R^n$

        \begin{itemize}
            \item naïve algo: brute force $O(n^4)$
            \item Warshall's: $w_{i, j} \gets w_{i, k} \land w_{k, j}$, $O(n^3)$ \\
                interior vertices: exclude $a, b$ \\
                init: $W = M_R$
        \end{itemize}

    \section{Equivalence Relations}
    \subsection{definitions}
        \textbf{Equivalence Relation}: a relation $R$ on $A$ that is:\\
        a. reflexive: $a = a$\\
        b. symmetric: $a = b \land b = a$\\
        c. transitive: $a = b \land b = c \implies a = c$

        % Examples: congruence / similarity of triangles; $a \equiv b \pmod n$ on $\field{Z}$; people having the same birthday.

        \textbf{Equivalence Class}: $[a] = R(a) = \{x \in A : a R x\}$. \\
        Every $x \in [a]$ is a \textbf{representative} of the class, and $[a] = [b]$ whenever $b \in [a]$.

        \textbf{Quotient Set}: $A/R = \{[a] : a \in A\}$, the set of all distinct equivalence classes.

        \textbf{Lemma}. $a R b \iff R(a) = R(b)$.

        \textbf{Theorem}. Let $R$ be an equivalence relation on $A$. Then $A/R$ is a partition of $A$, and $R$ is exactly the relation ``lying in the same block''. \\
        Conversely, every partition $P$ of $A$ determines an equivalence relation by
        $$ a R b \iff a \text{ and } b \text{ are in the same block of } P. $$
        Hence equivalence relations on $A$ and partitions of $A$ are in one-to-one correspondence.（集合划分和等价关系一一对应）

        \subsection{Congruence Classes}
            $$ a \equiv b \pmod n \iff n \mid (a - b) $$
            It is an equivalence relation on $\field{Z}$, and moreover a congruence relation on the semigroup $(\field{Z}, +)$.

            The classes are $[0], [1], \cdots, [n-1]$. The quotient set is denoted $\field{Z}_n$ (also written $\field{Z}/n\field{Z}$), and
            $$ [a] \oplus [b] = [a + b] $$

        \subsection{Operations on Equivalence Relations}
            Let $R$ and $S$ be equivalence relations on $A$.
            \begin{itemize}
                \item $R \cap S$ is an equivalence relation;
                \item $R \cup S$ don't need to be transitive; the smallest equivalence relation containing both is the transitive closure $t(R \cup S)$;
                \item $R^{-1} = R$ and $R \circ R = R$ (symmetry + transitivity).
            \end{itemize}

    \section{Partial Orderings}
    \subsection{definitions}
        \textbf{Partial Order}: a relation $R$ on $A$ that is:\\
        a. reflexive: $a \leqslant a$\\
        b. antisymmetric: $a \leqslant b \land b \leqslant a \implies a = b$\\
        c. transitive: $a \leqslant b \land b \leqslant c \implies a \leqslant c$

        The pair $(A, R)$ is a \textbf{partially ordered set}, or \textbf{poset}; usually $R$ is written $\leqslant$.

        % Examples: $(\mathcal{P}(S), \subseteq)$; $(\field{Z}^+, \mid)$ (divisibility); $(\field{R}, \le)$.

        \textbf{Dual Poset}: $(A, \geqslant)$, where $a \geqslant b \iff b \leqslant a$. The dual of a poset is again a poset, and its Hasse diagram is the original one turned upside down.

        \textbf{Comparable}: $a$ and $b$ are comparable if $a \leqslant b$ or $b \leqslant a$; otherwise they are \textbf{incomparable}. \\
        \textbf{Chain}: a poset in which every two elements are comparable; the order is a \textbf{linear} (or \textbf{total}) order. 称为线序或全序。

        \textbf{Product Order}: on $A \times B$,
        $$ (a, b) \leqslant (a', b') \iff a \leqslant a' \text{ in } A \text{ and } b \leqslant b' \text{ in } B. $$
        \textbf{Lexicographic Order}: $(a, b) \prec (a', b') \iff a < a'$, or $a = a'$ and $b \leqslant b'$. \\
        On $S^n$ this is exactly the dictionary order on words of length $n$.

        \textbf{Well-Ordered Set}: a poset in which every nonempty subset has a least element. \\
        Every finite chain is well-ordered; $(\field{N}, \le)$ is well-ordered, but $(\field{Z}, \le)$ is not (because, for instance, $\field{Z}$ itself does not have a least element !).

        \textbf{Quasi-Order}: $R$ is irreflexive and transitive. 拟序是小于号 $<$，补上 $I_A$ 并计算闭包成为偏序。

        \subsection{Hasse Diagram}
            The \textbf{Hasse diagram} is a simplified digraph of a finite poset:
            \begin{itemize}
                \item delete every self-loop (reflexivity);
                \item delete every edge implied by transitivity: if $a < b < c$, delete the edge $a \to c$;
                \item draw all remaining edges upward and omit the arrows;
                \item replace the circles by dots.
            \end{itemize}
            Thus $a$ is drawn below $b$ and joined to it iff $a < b$ and no $c$ satisfies $a < c < b$; one says that $b$ \textbf{covers} $a$. (\textbf{Covering Relation})\\
            A finite chain is a vertical path.

            \begin{center}
            \begin{tikzpicture}[scale=0.85]
                \node (n12) at (0, 3) {$12$};
                \node (n4)  at (-1.5, 1.6) {$4$};
                \node (n3)  at (1.5, 1.6) {$3$};
                \node (n2)  at (-1.5, 0) {$2$};
                \node (n1)  at (0, 0) {$1$};
                \draw (n1) -- (n2);
                \draw (n1) -- (n3);
                \draw (n2) -- (n4);
                \draw (n3) -- (n12);
                \draw (n4) -- (n12);
            \end{tikzpicture}
            \end{center}
            e.g. the Hasse diagram of $(\{1, 2, 3, 4, 12\}, \mid)$.

        \subsection{Topological Sorting}
            A \textbf{topological sorting} of a finite poset $(A, \le)$ is a linear order $\prec$ on $A$ extending $\le$:
            $$ a \leqslant b \implies a \prec b $$
            Algorithm: Each time choose a minimal element of $A$, append it to the list and delete it from $A$.
            The result depends on the choices, but every finite poset has at least one topological sorting.

        \subsection{Extremal Elements}
            Let $(A, \le)$ be a poset and $B \subseteq A$.
            \begin{itemize}
                \item \textbf{maximal} / \textbf{minimal} 极大元/极小元
                \item \textbf{greatest} / \textbf{least} 最大元(unit element, $1$) / 最小元 (zero element, $0$)\\
                    仅对 $A$ 有定义。
                \item \textbf{upper bound} / \textbf{lower bound} 上/下界
                \item \textbf{LUB} / \textbf{GLB}: least upper bound (join) / greatest lower bound (meet) \\
                    $\text{lub} B \sim \sup B$; $\text{glb} B \sim \inf B$
            \end{itemize}

            Notes:
            \begin{itemize}
                \item greatest / least elements are unique when they exist; \\
                    If $x, y$ are both g/l, then by antisymmetric $x = y$. \\
                    $\field{Z}$ has neither of them.
                \item every finite nonempty poset has \textbf{at least one} maximal and one minimal element, but they need not be unique; \\
                    Maintain a chain of elements. If $a_1/a_n$ is not the $\min/\max$, then we can add elements to the front/tail. However, there are only finite elements, thus the process will end in finite steps.
                \item a maximal/minimal element need not be an upper/lower bound of the whole poset;
                \item $B$ may have many ub/lb but \textbf{at most one} LUB/GLB, and the LUB/GLB need not belong to $B$;\\
                    also by antisymmetric.
                \item in a Hasse diagram, the LUB of $B$ is the first vertex reachable from every $b \in B$ by upward paths; the GLB is found by downward paths.
            \end{itemize}

        \subsection{Lattices}
            A \textbf{lattice} is a poset $(L, \le)$ in which every pair $\{a, b\}$ has a LUB and a GLB:
            $$ a \vee b := \operatorname{LUB}\{a, b\} \quad (\text{join}), \qquad
               a \wedge b := \operatorname{GLB}\{a, b\} \quad (\text{meet}). $$
            Thus a lattice is an algebraic structure with two binary operations.

            Examples: $(\mathcal{P}(S), \subseteq)$ with $\vee = \cup$, $\wedge = \cap$; \\
            $(\field{Z}^+, \mid)$ with $\vee = \operatorname{lcm}$, $\wedge = \gcd$; \\
            $(D_n, \mid)$, where $D_n$ is the set of positive divisors of $n$;\\
            every chain.

            \textbf{Sublattice}: a nonempty $S \subseteq L$ closed under $\vee$ and $\wedge$.

            \textbf{Properties}. For all $a, b, c \in L$:
            \begin{itemize}
                \item $a \vee b = b \iff a \leqslant b \iff a \wedge b = a$;
                \item idempotent: $a \vee a = a = a \wedge a$;
                \item commutative: $a \vee b = b \vee a$, \ $a \wedge b = b \wedge a$;
                \item associative: $a \vee (b \vee c) = (a \vee b) \vee c$, and dually;
                \item absorption: $a \vee (a \wedge b) = a$, \ $a \wedge (a \vee b) = a$;
                \item order-preserving: $a \leqslant b \implies a \vee c \leqslant b \vee c$ and $a \wedge c \leqslant b \wedge c$.
            \end{itemize}

            To prove: use definition (e.g. $a \leqslant (a \vee b)$), and reflexivity ($a \leqslant a$). prove by contradiction.

            Special lattices:
            \begin{itemize}
                \item \textbf{bounded}: has a least element $0$ and a greatest element $1$;
                \item \textbf{distributive}: $a \wedge (b \vee c) = (a \wedge b) \vee (a \wedge c)$, and the dual law;
                \item \textbf{complemented}: bounded, and every $a$ has a complement $a'$ with $a \vee a' = 1$, $a \wedge a' = 0$.
            \end{itemize}

        \subsection{Finite Boolean Algebras}
            $B_n$ is the lattice of all $n$-tuples of $0$'s and $1$'s under the componentwise order; equivalently $B_n \cong (\mathcal{P}(S), \subseteq)$ with $|S| = n$, so $|B_n| = 2^n$.
            $$ x \leqslant y \iff x_k \leqslant y_k \; \forall k, \qquad
               x \wedge y = \min, \qquad x \vee y = \max, \qquad x' = 1 - x. $$

            A finite lattice is a \textbf{Boolean algebra} if it is isomorphic to $B_n$ for some $n$. \\
            $D_n$ is a Boolean algebra $\iff$ $n$ is a product of distinct primes (no prime square divides $n$).\\
            ($D_n$ is the set of positive divisors of $n$)

            A Boolean algebra satisfies the usual set identities with $\wedge, \vee, {}'$ in place of $\cap, \cup, \overline{\phantom{A}}$: De Morgan's laws, involution $(x')' = x$, distributivity, $x \vee 0 = x$, $x \wedge 1 = x$, etc.

\chapter{Semigroups and Groups}

    \section{Binary Operations}
        A \textbf{binary operation} on a nonempty set $A$ is an everywhere-defined function $*: A \times A \to A$; one writes $a * b$ instead of $*(a, b)$. \\
        A unary operation is a function $A \to A$; an $n$-ary operation is a function $A^n \to A$.
        
        For a finite set $A = \{a_1, \cdots, a_n\}$ an operation is given by its \textbf{operation table}; \\
        There are $|A|^{|A|^2}$ binary operations on a set $A$.
        
        Laws of binary operations:
        \begin{itemize}
            \item \textbf{commutative}: $a * b = b * a$; \\
                iff the table is symmetric about the main diagonal.
            \item \textbf{associative}: $(a * b) * c = a * (b * c)$;
            \item \textbf{idempotent}: $a * a = a$; \\
                主对角线上的每一个元素都与它所在行（列）的表头元素相同。
            \item \textbf{distributive} over $+$: $a * (b + c) = (a * b) + (a * c)$ (and dually on the right); \\
                左分配（第一分配律），右分配（第二分配律）
            \item \textbf{cancellation}: $a * b = a * c \implies b = c$; $b * a = c * a \implies b = c$ \\
                左消去律，右消去律
            \item \textbf{absorption}: $a * (a + b) = a$.
        \end{itemize}

        \textbf{Identity} $e$: $e * a = a * e = a$ for all $a$. \\
        It is unique. $e_1 = e_1 * e_2 = e_2$

        \textbf{Zero} $z$: $z * a = a * z = z$ for all $a$

        \textbf{Idempotent Element}: $a * a = a$

        \textbf{Inverse} of $a$: an element $a'$ with $a * a' = a' * a = e$. \\
        It is unique, if associativity is satisfied. $(x * e) * y = e * y = y, x * (e * y) = x * e = x$, thus $x = y$.

    \section{Semigroups and Monoids}
        A \textbf{semigroup} 半群 $(S, *)$ is a nonempty set together with an associative binary operation; it is \textbf{commutative} if $*$ is commutative. \\
        Examples: $(\field{Z}, +)$, $(\field{Z}^+, \cdot)$, $(\mathcal{P}(S), \cup)$, $(S^S, \circ)$, and the free semigroup $(A^*, \cdot)$ of words over an alphabet $A$ under catenation.

        Because of associativity, all products $a_1 * a_2 * \cdots * a_n$ are equal regardless of parentheses.

        \textbf{Monoid} 幺半群: a semigroup with an identity $e$. \\
        Examples: $(\field{N}, +)$ with $e = 0$; $(\field{Z}, \cdot)$ with $e = 1$; $(\mathcal{P}(S), \cup)$ with $e = \emptyset$; $(A^*, \cdot)$ with $e = \varepsilon$ (the empty word).

        \textbf{Subsemigroup}: a nonempty $T \subseteq S$ closed under $*$. \\
        \textbf{Submonoid}: a subsemigroup $T$ that also contains $e$.

        Powers: $a^1 = a$, $a^n = a^{n-1} * a$; in a monoid $a^0 = e$, and $a^m * a^n = a^{m+n}$. \\
        The set $\langle a \rangle = \{a^i : i \geqslant 0\}$ is the submonoid generated by $a$.

        \textbf{Isomorphism / Homomorphism}: 同构 / 同态\\
        for semigroups $(S, *)$ and $(T, *')$, a map $f: S \to T$ is a homomorphism if
        $$ f(a * b) = f(a) *' f(b) \quad \forall a, b \in S; $$
        it is an isomorphism if it is also bijective. Then $S \cong T$.
        An isomorphism preserves the identity (if it exists), commutativity, and subsemigroups. \\
        An onto homomorphism also preserves monoids, commutativity, and subsemigroups; it need not be injective.

    \section{Groups}
        A \textbf{group} $(G, *)$ is a monoid in which every element has an inverse:
        \begin{enumerate}
            \item $(a * b) * c = a * (b * c)$;
            \item there is $e \in G$ with $a * e = e * a = a$ for all $a$;
            \item for every $a \in G$ there is $a^{-1} \in G$ with $a * a^{-1} = a^{-1} * a = e$.
        \end{enumerate}
        $G$ is \textbf{Abelian} if $a * b = b * a$ for all $a, b$. We usually write $ab$ for $a * b$.

        Examples: $(\field{Z}, +)$, $(\field{Q} \setminus \{0\}, \cdot)$, $(\field{R} \setminus \{0\}, \cdot)$, $(\field{Z}_n, \oplus)$; \\
        the symmetric group $(S_n, \circ)$ of all permutations of $n$ letters (order $n!$); \\
        the dihedral group $(D_n, \circ)$ of symmetries of a regular $n$-gon; \\
        the Klein four-group $V$, and the cyclic group $\field{Z}_4$.

        \textbf{Elementary properties}:
        \begin{itemize}
            \item the inverse is unique, $(a^{-1})^{-1} = a$, and $(ab)^{-1} = b^{-1} a^{-1}$;
            \item cancellation: $ab = ac \implies b = c$ and $ba = ca \implies b = c$;
            \item $ax = b$ has the unique solution $x = a^{-1}b$; $ya = b$ has the unique solution $y = ba^{-1}$;
            \item in the multiplication table of a finite group, every element appears exactly once in each row and each column.
        \end{itemize}

        \textbf{Order}: $|G|$ is the number of elements of $G$. \\
        The \textbf{order of $a \in G$} is the least $n \geqslant 1$ with $a^n = e$ (if no such $n$ exists, the order is infinite). In a finite group every element has finite order.

        Up to isomorphism there is exactly one group of each order $1, 2, 3$, two of order $4$ (namely $V$ and $\field{Z}_4$), and two of order $6$ (namely $\field{Z}_6$ and $S_3$).

    \section{Subgroups and Cosets}
        A \textbf{subgroup} of $G$ is a subset $H \subseteq G$ such that
        \begin{enumerate}
            \item $e \in H$;
            \item $a, b \in H \implies ab \in H$;
            \item $a \in H \implies a^{-1} \in H$.
        \end{enumerate}
        Equivalently, $H \neq \emptyset$ and $a, b \in H \implies ab^{-1} \in H$; then $(H, *)$ is itself a group. Write $H \leqslant G$.

        Examples: $G$ and $\{e\}$ (the \textbf{trivial subgroups}); $n\field{Z} \leqslant \field{Z}$; $A_n \leqslant S_n$ (alternating group); the cyclic subgroup $\langle a \rangle = \{a^i : i \in \field{Z}\}$.

        \textbf{Coset}: for $H \leqslant G$ and $a \in G$,
        $$ aH = \{ah : h \in H\} \quad (\text{left coset}), \qquad
           Ha = \{ha : h \in H\} \quad (\text{right coset}). $$
        The left cosets of $H$ form a partition of $G$, and $|aH| = |H|$ for every $a$.

        \textbf{Lagrange's Theorem}. If $G$ is finite and $H \leqslant G$, then $|H|$ divides $|G|$, and the number of distinct left cosets, the \textbf{index} $[G : H]$, equals $|G| / |H|$. \\
        Corollaries: the order of every element divides $|G|$; a group of prime order is cyclic.

        \textbf{Normal subgroup}: $H \trianglelefteq G$ if $aH = Ha$ for all $a \in G$, equivalently $a^{-1}Ha = H$. Every subgroup of an Abelian group is normal.

        \textbf{Quotient group}: if $N \trianglelefteq G$, then
        $$ G/N = \{aN : a \in G\}, \qquad (aN)(bN) = abN $$
        is a group, with identity $N$ and inverse $(aN)^{-1} = a^{-1}N$.

    \section{Homomorphisms and Quotient Groups}
        For groups $(G, *)$ and $(H, \circ)$, a map $f: G \to H$ is a \textbf{homomorphism} if
        $$ f(a * b) = f(a) \circ f(b) \quad \forall a, b \in G. $$
        A bijective homomorphism is an \textbf{isomorphism}, written $G \cong H$.

        Basic facts: $f(e_G) = e_H$, $f(a^{-1}) = f(a)^{-1}$, $f(G) \leqslant H$, and $f$ is injective $\iff \ker f = \{e_G\}$.

        \textbf{Kernel}: $\ker f = \{a \in G : f(a) = e_H\}$; it is a normal subgroup of $G$.

        \textbf{Fundamental Homomorphism Theorem}. If $f: G \to H$ is onto, then
        $$ G / \ker f \cong H, \qquad a \ker f \mapsto f(a). $$
        The map $a \mapsto aN$ is the \textbf{natural homomorphism} $G \to G/N$.

        More generally, for a semigroup $S$ and a \textbf{congruence relation} $R$ on $S$ (an equivalence relation compatible with $*$: $a R a'$ and $b R b'$ imply $ab R a'b'$), the quotient $S/R$ is a semigroup under $[a] * [b] = [ab]$, and $a \mapsto [a]$ is an onto homomorphism. \\
        For groups, congruence relations are exactly the relations
        $$ a R b \iff a^{-1}b \in N $$
        for some normal subgroup $N$.

        \textbf{Products}: $G_1 \times G_2$ with componentwise operation is a group; $G_1 \times G_2 \cong G_2 \times G_1$, and
        $$ \field{Z}_m \times \field{Z}_n \cong \field{Z}_{mn} \iff \gcd(m, n) = 1. $$

    \section{Rings and Fields}
        A \textbf{ring} $(R, +, \cdot)$ is a set with two binary operations such that $(R, +)$ is an Abelian group, $\cdot$ is associative, and $\cdot$ distributes over $+$. \\
        It is \textbf{commutative} if $\cdot$ is commutative, and a \textbf{ring with identity} if $(R, \cdot)$ is a monoid.

        Examples: $(\field{Z}, +, \cdot)$ (commutative, with identity); $(\field{Z}_n, +, \cdot)$; \\
        $(M_n(\field{R}), +, \cdot)$ (not commutative); $(\field{Q}, +, \cdot)$.

        \textbf{Field}: a commutative ring with identity in which every nonzero element has a multiplicative inverse. \\
        Examples: $\field{Q}$, $\field{R}$, $\field{C}$; $\field{Z}_p$ is a field $\iff$ $p$ is prime.

        \begin{longtable}[h]{l l}
            \toprule
            Structure & Axioms \\
            \midrule
            Semigroup & associative binary operation \\
            Monoid & semigroup + identity \\
            Group & monoid + inverses \\
            Abelian group & group + commutative operation \\
            Ring & Abelian group under $+$, associative $\cdot$, distributive \\
            Field & commutative ring with identity, nonzero elements invertible \\
            \bottomrule
        \end{longtable}

\end{document}